\documentclass[11pt]{article}
\usepackage{preamble}
\setlength{\parskip}{4pt}

\begin{document}

\daybanner{AP Statistics \textperiodcentered\ Unit 4 \textperiodcentered\ Topic 4.2 \textperiodcentered\ B: 2027-01-12 \textperiodcentered\ E: 2027-02-26 \textperiodcentered\ TEACHER KEY}{Can Chance Explain the Mean Difference?}

\sectionbanner{CED TETHER}

{\small
\begin{itemize}[leftmargin=1.5em,itemsep=2pt,topsep=2pt]
  \item \textbf{VAR-1.I} Identify questions suggested by probabilities of errors in statistical inference. \textbf{[Skill 1.A]}
  \item \textbf{VAR-1.I.1} Random variation may result in errors in statistical inference.
\end{itemize}
}
\textbf{Essential question:} How can rerandomization distinguish a treatment explanation from chance variation while preserving the possibility of inferential error?\par
\textbf{DOK-3 skill:} 4.B \textperiodcentered\ \textbf{FRQ pattern:} {\small\texttt{rerandomization-\allowbreak{}probability-\allowbreak{}error-\allowbreak{}and-\allowbreak{}scope}}\par
\textbf{DOK rationale:} Part (c) adjudicates competing causal claims by coordinating the observed mean difference, a rerandomization tail proportion, its conditional meaning, the remaining chance-error risk, and the distinct roles of random assignment and volunteer recruitment.\par

\frameworkphaseheader{Do Now (first take) $\rightarrow$ Explore (video + follow-along) $\rightarrow$ Exit (finish + turn in)}{1 $\rightarrow$ 3}{45}{%
\begin{itemize}[leftmargin=1.5em,itemsep=2pt,topsep=2pt]
  \item Display the board and ask for one sentence using the study description.
  \item Begin the 7.1 video follow-along at minute 5.
  \item Return to the board at minute 33. Students finish the shortened parts and justify their judgment.
  \item Collect at the door. Score part (c) only using the three required elements.
\end{itemize}
}{%
\begin{itemize}[leftmargin=1.5em,itemsep=2pt,topsep=2pt]
  \item Write a one-sentence first take before the video.
  \item Complete the video follow-along worksheet.
  \item Finish (a)--(c), revise the first take in the exit line, and turn in the sheet.
\end{itemize}
}{%
\begin{itemize}[leftmargin=1.5em,itemsep=2pt,topsep=2pt]
  \item Which specific number or study detail supports your choice?
  \item What claim would go beyond the evidence, and why?
\end{itemize}
}{Circulate and ask for evidence. Accept a concise response; do not require omitted calculations or a second full inference write-up.}

\begin{callout}{\IconWarn\ WATCH FOR}{calloutred}
\begin{itemize}[leftmargin=1.5em,itemsep=2pt,topsep=2pt]
  \item Choosing a speaker without a reason.
  \item Copying numbers without explaining what they support.
\end{itemize}
\end{callout}

\sectionbanner{SCORING PART (c) --- E / P / I}

\textbf{Expected elements}
\begin{itemize}[leftmargin=1.5em,itemsep=2pt,topsep=2pt]
  \item \textbf{uses-simulation} (required) --- Supports Nia using 18/2,000 under no effect
  \item \textbf{allows-error} (required) --- Explains that rare under no effect does not make an error impossible
  \item \textbf{limits-scope} (required) --- Limits the causal claim to the volunteers rather than all district students
\end{itemize}

{\small\renewcommand{\arraystretch}{1.25}
\noindent\begin{tabular}{|p{0.5in}|p{5.9in}|}
\hline
\textbf{E} & Uses the conditional simulation evidence, allows an error, and limits the causal claim to the volunteers. \\ \hline
\textbf{P} & Shows substantive valid reasoning for at least one required element, but does not meet all E requirements. Credit correct reasoning even when another claim is wrong. \\ \hline
\textbf{I} & Shows no substantive valid reasoning for any required element; a name, unsupported choice, or copied number alone does not count. \\ \hline
\end{tabular}}

\clearpage
\focusbanner{TODAY'S PROBLEM --- annotated}

Suppose 50 student volunteers test a planning prompt for a logic task. They are randomly assigned to two groups of 25. Mean times are 14.2 minutes with the prompt and 17.8 with standard directions: prompt minus standard is $-3.6$. Under no effect, a simulation shuffles the same times into new groups. Only 18 of 2{,}000 shuffled differences are $-3.6$ or lower.\par\smallskip \textbf{Nia:} ``This is evidence that the prompt lowered mean time for these volunteers, but we could still be wrong.''\par \textbf{Owen:} ``The $0.009$ is the probability of no effect, so there is a $99.1\%$ chance the prompt works for every district student.''\par
\begin{center}
\textbf{Experiment results}\par\smallskip
\begin{tabular}{|l|c|c|}
\hline
\textbf{Group} & \textbf{n} & \textbf{Mean (min)} \\ \hline
\textbf{Prompt} & 25 & 14.2 \\ \hline
\textbf{Standard} & 25 & 17.8 \\ \hline
\end{tabular}
\end{center}
\begin{firsttakebox}{FIRST TAKE (5 min)}
Whose claim sounds more believable, Nia's or Owen's? Give one reason from the study.\par
\answer{Ungraded. Everyday language is enough before the video; formal vocabulary belongs in the finish.}
\end{firsttakebox}

\textit{Rules callout ``ASK WHAT CHANCE WOULD DO UNDER NO EFFECT (keep on the page)'' is printed on the student sheet.}\par\medskip

\dokbadge{1}~\textbf{(a)} What difference in group means would we expect if the prompt had no effect?\par
\answer{The expected prompt-minus-standard difference is 0 minutes.}

\dokbadge{2}~\textbf{(b)} Compute $18/2{,}000$. In one sentence, explain what this proportion estimates if the prompt has no effect.\par
\answer{$18/2{,}000=0.009=0.9\%$. If the prompt has no effect, about $0.9\%$ of random assignments give a difference of $-3.6$ minutes or lower.}

\dokbadge{3}~\textbf{(c)} In 3--4 sentences, decide whose claim is better supported. Use the simulation result, explain why an error is still possible, and say why the volunteers do not represent every district student.\par
\answer{Nia is better supported: a difference of $-3.6$ or lower occurred in only 18 of 2{,}000 no-effect trials. This estimates the chance of such a result assuming no effect, not the probability that no effect is true, so an incorrect effect claim remains possible. Random assignment supports a causal conclusion for these volunteers. Because they volunteered, the study does not justify a claim about every district student.}

\begin{summaryexitbox}{EXIT}
One thing the video changed about my first take: \hrulefill
\end{summaryexitbox}

\end{document}
